\documentclass[11pt]{article}
\usepackage{fontspec}
\usepackage[french,english]{babel}
\usepackage{amsmath}
\usepackage{amssymb}
\usepackage{amsthm}
\usepackage{mathrsfs}
\usepackage{tikz-cd}
\usepackage{enumitem}
\usetikzlibrary{babel}

\newcommand{\printedpage}[1]{\marginpar{\footnotesize p.~#1}}
\newcommand{\msp}[1]{\marginpar{\footnotesize [#1]}}

\newcommand{\HOT}{\mathrm{HOT}}
\newcommand{\Cat}{\mathrm{Cat}}
\newcommand{\Hot}{\mathrm{Hot}}
\newcommand{\hot}{\mathrm{hot}}
\newcommand{\Ob}{\operatorname{Ob}}
\newcommand{\Hom}{\operatorname{Hom}}
\newcommand{\ul}[1]{\underline{\mathrm{#1}}}
\newcommand{\cercle}[1]{\tikz[baseline=(c.base)]\node[draw,circle,inner sep=1pt](c){#1};}
\newcommand{\ovale}[1]{\tikz[baseline=(c.base)]\node[draw,rounded corners=0.8ex,inner sep=2pt](c){#1};}

\begin{document}

\setcounter{footnote}{45}

\printedpage{44}
% head: A. Grothendieck : Dérivateurs, Ch. VI, Édition des Universités de Montpellier II et Paris VII
% foot: Version 9 novembre 2001
and the localizers $W_S^{\mathrm{g}}$ in these categories ($f\in W_S^{\mathrm{g}}$ if and only if the $f/s : X/s\longrightarrow Y/s$ are in $W$). We find a ``same localized category'', denoted
\[ \HOT^o_W(S), \]
and we have tautologically
\msp{page 71}
\[ \HOT^o_W(S)\approx\HOT_W(S^o). \]
\texttt{[rather} $(\HOT_W(S^o))^o$ \texttt{cf. p. 50]}

\ovale{III}\quad On the six inner categories of the diagram, the three notions of $W$-equivalence, in terms of $W_S$, $W_S^{\mathrm{d}}$, $W_S^{\mathrm{g}}$ coincide, and coincide with fiberwise $W$-equivalence $W_S^{\mathrm{f}}$. On the three categories of $(*)$ other than $\Cat/S$, the localizers $W_S^{\mathrm{f}}$, $W_S^{\mathrm{d}}$ coincide; on the three categories of $(*')$ other than $\Cat/S$, the localizers $W_S^{\mathrm{f}}$, $W_S^{\mathrm{g}}$ coincide. But in general (\footnote{probably in all cases where $S$ is not equivalent to a discrete category.% [sic] dicrète -> discrète, faute d'orthographe
}), on these six categories, these localizers $W_S^{\mathrm{f}}$ are not equal to $W_S$, more precisely give nonequivalent localized categories, they are therefore each equipped with \textit{two} essentially distinct localizers, $W_S^{\mathrm{f}}$ and $W_S$. The category $\Cat/S$ is the only category in the diagram on which the localizers $W_S$, $W_S^{\mathrm{d}}$ and $W_S^{\mathrm{g}}$ are all three essentially distinct. % [sic] essentiellment -> essentiellement, lettre omise
It is the only one that allows one to obtain the three localized categories simultaneously, as
\[
\begin{tikzcd}[column sep=small]
 & \HOT^{\mathrm{lc}}_W(S)\simeq(\Cat/S)W_S^{-1} & \\
\HOT_W(S)\simeq(\Cat/S)(W_S^{\mathrm{d}})^{-1} \arrow[ur,"\rho"] & & \HOT^o_W(S)\simeq(\Cat/S)(W_S^{\mathrm{g}})^{-1} \arrow[ul,"\sigma"']
\end{tikzcd}
\]

Here, $\rho$ and $\sigma$ are localization functors (as is evident from the definition).

Any functor or composite functor in the diagram on page 8, going from an inner category of the diagram to a bordering category, induces a functor on the localized categories for $W_S^{\mathrm{d}}$ (when the target category is one of the four on the left,
\msp{page 72}
recalled in $(*)$, p. 70), or for $W_S^{\mathrm{g}}$ (when these categories are among the four on the right-hand side, recalled in $(*')$), which localizers also coincide with $W_S^{\mathrm{f}}$ except in the sole case of the category $\Cat/S$. All the functors thus obtained reduce, up to equivalence, to two, according as the target localized category is $\HOT_W(S)$ (case of localization $W_S^{\mathrm{d}}$), or is $\HOT^o_W(S)$ (case of localization $W_S^{\mathrm{g}}$). We thus find functors (in the direction opposite to $\rho$ and to $\sigma$)
\printedpage{45}
% head: A. Grothendieck : Dérivateurs, Ch. VI, Édition des Universités de Montpellier II et Paris VII
% foot: Version 9 novembre 2001
\[
\begin{tikzcd}
 & \HOT^{\mathrm{lc}}_W(S) \arrow[dl,"i"'] \arrow[dr,"j"] & \\
\HOT_W(S) & & \HOT^o_W(S).
\end{tikzcd}
\]

The functors $\rho$, $\sigma$ are interchanged by duality $S\longmapsto S^o$, and likewise for $i$, $j$.

\ovale{IV}\quad The functors $\rho$, $i$ are adjoint to one another in that order, i.e. we have for two objects $\xi\in\Ob\HOT_W(S)$, $\eta\in\Ob\HOT^{\mathrm{lc}}_W(S)$
\[ \Hom(\rho(\xi),\eta)\simeq\Hom(\xi,i(\eta)) \]
(isomorphism of bifunctors). And dually the functors $\sigma$, $j$ are adjoint to one another % [sic] adjoint -> adjoints, accord du pluriel
in that order: for $\xi\in\Ob\HOT^o_W(S)$, $\eta\in\Ob\HOT^{\mathrm{lc}}_W(S)$
\[ \Hom(\xi,j(\eta))\simeq\Hom(\sigma(\xi),\eta). \]

\msp{page 73}
It follows in particular that \textit{the functors $i$, $j$ are fully faithful.}

\cercle{V}\quad We have obvious functors
\[
\begin{aligned}
&(*) &\qquad& \HOT_W(S) &\quad& \longrightarrow &\quad& \underline{\mathrm{Hom}}(S^o,\Hot_W)\\
&(*') && \HOT^o_W(S) && \longrightarrow && \underline{\mathrm{Hom}}(S,\Hot_W).
\end{aligned}
\]
The first is interpreted on the four categories $(*)$ (bordering on the left) on page 70, which serve indifferently to describe $\HOT_W(S)$, by % [sic] indifférement -> indifféremment, faute d'orthographe
\[ s\longmapsto\hot_W(X_s) \]
in the case of the three categories distinct from $\Cat/S$, and by
\[ s\longmapsto\hot_W(s\backslash X) \]
in the case of $\Cat/S$. The functorial dependence on $s$ is obvious in the case % [sic] dépendence -> dépendance, faute d'orthographe
of $\Cat/S$ (given by $s\longmapsto\hot_W(s\backslash X)$), and in the case of $\ul{Fibsc}\,S$, $\ul{Fib}\,S$, where we have, for $u : s\longrightarrow t$ in $S$, the base change
\[ u^* : X_t\longrightarrow X_s. \]
In the case of $\ul{Liss}\,S$, the functorial dependence is deduced from that of $\hot_W(s\backslash X)$, thanks % [sic] dépendence -> dépendance, faute d'orthographe
to the fact that the inclusion
\[ X_s\longrightarrow s\backslash X \]
is then a $W$-equivalence.

We have a dual interpretation of the map
\printedpage{46}
% head: A. Grothendieck : Dérivateurs, Ch. VI, Édition des Universités de Montpellier II et Paris VII
% foot: Version 9 novembre 2001
\msp{page 74}
$(*)'$, \texttt{[illegible word]}
\[ s\longmapsto\hot_W(X_s)\qquad\text{or}\qquad s\longmapsto\hot_W(s\backslash X), \]
with the functorial law in $s$ deduced from the functor $t\backslash X\longrightarrow s\backslash X$ \texttt{[rather} $X/s\longrightarrow X/t$\texttt{]} (for $u : s\longrightarrow t$) in the case of $\Cat/S$, from the cobase change functor in the cases of $\ul{Cofib}\,S$ and $\ul{Cofibsc}\,S$, and from the variant for $\Cat/S$ and from the $W$-equivalence
\[ X_s\longrightarrow X/s \]
in the case of $\ul{Prop}\,S$.

We then find this: for an object $\xi$ of $\HOT_W(S)$ (resp. $\HOT^o_W(S)$) to be in the essential image of $\HOT^{\mathrm{lc}}_W(S)$, it is necessary and sufficient that the associated functor $S^o\longrightarrow\Hot_W$ (resp. $S\longrightarrow\Hot_W$) be a local system (\footnote{the verification will still have to be written}), i.e. transform every arrow of $S$ into an isomorphism, i.e. factor through $(\Pi S)^o\longrightarrow\Hot_W$ resp. $\Pi S\longrightarrow\Hot_W$ (\footnote{where $\Pi S$ is the fundamental groupoid of $S$}). Thus for an $X$ in $\Cat/S$, this means either that the functors
\[ t\backslash X\longrightarrow s\backslash X \]
associated to $u : s\longrightarrow t$ in $S$ are all $W$-equivalences (case of $\HOT_W(S)$), or that the functors
\[ X/s\longrightarrow X/t \]
are (case of $\HOT^o_W(S)$). (NB. these conditions are by no means equivalent. But if $X$ \texttt{[is]} smooth \texttt{[over]} $S$,
\msp{page 75}
the first implies the second, and the converse is true if $X$ is proper over $S$, and in this case, the condition is equivalent to $X$ being in $\ul{Fibhot}$. Finally, if $X$ is already in $\ul{Fibhot}\,S$, then the two conditions are automatically satisfied (\footnote{We should also have introduced in the diagram $\ul{Liss}_0S$, $\ul{Prop}_0S$!}).
\[
\begin{tikzcd}[column sep=small, row sep=small]
\ul{Fibsc}_0S \arrow[dd,hook] \arrow[drr] & & & & & & & & \ul{Cofibsc}_0S \arrow[dll] \arrow[dd,hook] \\
 & & \ul{Fib}_0S \arrow[dd,hook] \arrow[dr,hook] & & \ul{Parf}\,S \arrow[dl,hook] \arrow[dd,hook] \arrow[dr,hook] & & \ul{Cofib}_0S \arrow[dl,hook] \arrow[dd,hook] & & \\
\ul{Fibsc}\,S \arrow[drr] & & & \ul{Liss}_0S \arrow[dr,hook] \arrow[dd,hook] & & \ul{Prop}_0S \arrow[dl,hook] \arrow[dd,hook] & & & \ul{Cofibsc}\,S \arrow[dll] \\
 & & \ul{Fib}\,S \arrow[dr,hook] & & \ul{Fibhot}\,S \arrow[dd,hook] \arrow[uuullll,dotted,bend right,"\Phi_0"'] \arrow[uuurrrr,dotted,bend left,"\Psi_0"'] & & \ul{Cofib}\,S \arrow[dl,hook] & & \\
 & & & \ul{Liss}\,S \arrow[dr,hook] & & \ul{Prop}\,S \arrow[dl,hook] & & & \\
 & & & & \Cat/S \arrow[uuullll,dotted,bend left,"\Phi"] \arrow[uuurrrr,dotted,bend right,"\Psi"'] & & & &
\end{tikzcd}
\]
\printedpage{47}
% head: A. Grothendieck : Dérivateurs, Ch. VI, Édition des Universités de Montpellier II et Paris VII
% foot: Version 9 novembre 2001
(\footnote{The diagram of \textit{fifteen} remarkable categories.}). Moreover, if two elements $\xi\in\Ob\HOT_W(S)$, $\xi'\in\Ob\HOT^o_W(S)$ come from the same $\eta\in\Ob\HOT^{\mathrm{lc}}_W(S)$, then the two corresponding functors
\[ (\Pi S)^o\xrightarrow{\ \varphi\ }\Hot_W\xleftarrow{\ \varphi'\ }\Pi S \]
are \textit{associated,} i.e. each is isomorphic to ``the inverse'' of the other.
\[
\left\{
\begin{array}{lcl}
\varphi^o(s) & = & \varphi(s)\\
\varphi^o(u) & = & \varphi(u)^{-1}.
\end{array}
\right.
\]

\ovale{VI}\quad NB In the new version of the diagram, where I have added the categories $\ul{Liss}_0S$ and $\ul{Prop}_0S$ to make the drawing more transparent, the categories called ``bordering'' become % [sic] appelés -> appelées, accord du participe
the ``ground floor'', and the ``inner'' categories, the ``first upper floor'' or ``upper floor''. All % [sic] rez-de-chausseée -> rez-de-chaussée, lettre en trop
7 squares of the diagram are cartesian. All the arrows are fully % [sic] catésiens -> cartésiens, lettre omise
faithful inclusions, except the four outer horizontal arrows from the categories of Cat-fibered or Cat-cofibered categories that are \textit{split}, with morphisms exclusively cartesian or cocartesian and even compatible with the splittings, to the corresponding categories of nonsplit structures.

\msp{page 76}
\[
\begin{tikzcd}[column sep=small, row sep=small]
 & & & & \ul{Fibhot}\,S \arrow[dd,hook,"\delta"] \arrow[dddllll,dotted,bend right,"\Phi_0"'] \arrow[dddrrrr,dotted,bend left,"\Psi_0"] & & & & \\
 & & & \ul{Liss}_0S \arrow[ur,hook] \arrow[dd,hook] & & \ul{Prop}_0S \arrow[ul,hook] \arrow[dd,hook] & & & \\
 & \ul{Fib}_0S \arrow[urr,hook] \arrow[dd,hook] \arrow[ddrrr,dotted,bend left] & & & \Cat/S \arrow[dd,"C"] \arrow[dddllll,dotted,bend right,"\Phi"'] \arrow[dddrrrr,dotted,bend left,"\Psi"] & & & \ul{Cofib}_0S \arrow[ull,hook] \arrow[dd,hook] \arrow[ddlll,dotted,bend right] & \\ % Reading verified against PDF leaf 9: sources Fib_0 S and Cofib_0 S, target middle Parf S.
\ul{Fibsc}_0S \arrow[ur] \arrow[dd,hook] & & & \ul{Liss}\,S \arrow[ur,hook] & & \ul{Prop}\,S \arrow[ul,hook] & & & \ul{Cofibsc}_0S \arrow[ul] \arrow[dd,hook] \\
 & \ul{Fib}\,S \arrow[urr,hook] & & & \ul{Parf}\,S \arrow[uuul] \arrow[uuur] \arrow[dd,dotted,"\mathrm{id}"] & & & \ul{Cofib}\,S \arrow[ull,hook] & \\
\ul{Fibsc}\,S \arrow[ur] & & & & & & & & \ul{Cofibsc}\,S \arrow[ul] \\
 & & & & \ul{Parf}\,S \arrow[uuul,dotted,hook] \arrow[uuur,dotted,hook] & & & &
\end{tikzcd}
\]

I have drawn in green dotted lines \texttt{[color in the original only]} the five functors $\Phi$, $\Psi$, $\Phi_0$, $\Psi_0$, $C$ which by passage to the localized categories give quasi-inverses for all the equivalences of categories stated in \cercle{I} \cercle{II} of the present Scholie. To be precise, the functor
\[ \Phi : \Cat/S\longrightarrow\ul{Fibsc}\,S \]
associates to every $X$ in $\Cat/S$ the split category over $S$
\[ s\longmapsto s\backslash X \]
\printedpage{48}
% head: A. Grothendieck : Dérivateurs, Ch. VI, Édition des Universités de Montpellier II et Paris VII
% foot: Version 9 novembre 2001
over $S$. This functor, and its restrictions to $\ul{Liss}\,S$, $\ul{Fib}\,S$ and, more generally, to any full subcategory $\mathcal{F}$ of $\Cat/S$ containing $\ul{Fib}\,S$, is compatible with the localizers $W_S$ on the one hand (on $\mathcal{F}$ and $\ul{Fibsc}\,S$), $W_S^{\mathrm{d}}$ on the other hand, and therefore passes to the corresponding localized categories, to provide the desired quasi-inverse of
\[ (\ul{Fibsc}\,S)\Sigma^{-1}\longrightarrow\mathcal{F}\Sigma^{-1}\qquad(\Sigma=W_S\text{ or }W_S^{\mathrm{d}}) \]

We can also, for any intermediate full subcategory $\mathcal{F}'$
\[ \ul{Fib}\,S\subseteq\mathcal{F}'\subseteq\mathcal{F}\subseteq\Cat/S, \]
regard $\Phi$ as a functor $\mathcal{F}\longrightarrow\mathcal{F}'$ (more precisely, take the composite
\[ \mathcal{F}\xrightarrow{\Phi|\mathcal{F}}\ul{Fibsc}\,S\longrightarrow\ul{Fib}\,S\hookrightarrow\ \mathcal{F}'\ ), \]
say $\Phi_{\mathcal{F}',\mathcal{F}}$, and from what has just been said it follows that the two functors
\[
\begin{tikzcd}
\mathcal{F}' \arrow[r,hook] & \mathcal{F} \arrow[l,bend left,"\Phi_{\mathcal{F}',\mathcal{F}}"]
\end{tikzcd}
\]
\msp{page 77}
establish equivalences quasi-inverse to one another
\[
\begin{tikzcd}
\mathcal{F}'\Sigma^{-1} \arrow[r,"\approx"] & \mathcal{F}\Sigma^{-1} \arrow[l,bend left,"\approx"',"\Phi_{\mathcal{F}',\mathcal{F}}"]
\end{tikzcd}
\qquad\qquad(\Sigma=W_S\text{ or }W_S^{\mathrm{d}})
\]

Dual statement for the functor
\[ \Psi : \Cat/S\longrightarrow\ul{Cofibsc}\,S, \]
associating to $X$ in $\Cat/S$ the split cofibered category over $S$
\[ s\longmapsto X/s. \]
It gives pairs of quasi-inverse functors
\[
\begin{tikzcd}
(\ul{Cofibsc}\,S)\Sigma^{-1} \arrow[r] & \mathcal{F}\Sigma^{-1} \arrow[l,bend left,"\Psi_{\mathcal{F}}"]
\end{tikzcd}
\qquad\qquad(\Sigma=W_S\text{ or }W_S^{\mathrm{g}})
\]
for any full subcategory $\mathcal{F}$ of $\Cat/S$ such that % [sic] tout -> toute, accord avec sous-catégorie
\[ \ul{Cofib}\,S\subseteq\mathcal{F}\subseteq\Cat/S, \]
and pairs of quasi-inverse functors

\end{document}
